% Response to reviewers — P1 V8 revision (JCIM)
% NOTE FOR AUTHORS: all reviewer-requested computations are complete as of 2026-09-09,
% including the R2.4 two-arm MTX-retained vs MTX-stripped DEKOIS re-run (honest-negative,
% integrated into the SM table and main text). The only remaining item is the R2.6
% repository release + Zenodo DOI (author action; checklist in P1_R26_RELEASE_CHECKLIST.md).

\documentclass[11pt]{article}
\usepackage[margin=2.5cm]{geometry}
\usepackage[T1]{fontenc}
\usepackage[utf8]{inputenc}
\usepackage{lmodern}
\usepackage{enumitem}
\usepackage{url}
\usepackage{xr}
\usepackage[colorlinks=true,linkcolor=blue,citecolor=blue,urlcolor=blue]{hyperref}
\usepackage[noabbrev,capitalize]{cleveref}
\externaldocument[M-]{P1_V8_main}
\externaldocument[SM-]{P1_V8_SM}
\setlength{\parskip}{4pt}

\title{\bfseries Response to Reviewers\\[4pt]
\large Revision of: ``Target breadth and mutation resilience in
African-natural-product-inspired antimalarial chemotypes: a computational analysis''
(submitted to \emph{Journal of Chemical Information and Modeling})}
\author{}
\date{}

\begin{document}
\maketitle

\noindent Dear Editor,

We thank both reviewers for their careful, constructive reading of our manuscript.
Reviewer~1 raised two major concerns and Reviewer~2 raised nine major points and seven
minor points. We have revised the manuscript (V8) to address every point. Below we
respond point by point; locations refer to the revised main text (M) or Supporting
Information (SI).

\section*{Reviewer 1}

\subsection*{R1.1 --- Uneven structural evidence quality across the four targets}
\textbf{Reviewer:} ``The four targets have vastly different levels of structural validation.
These differences would influence the `target breadth' result.''

\textbf{Response:} We agree, and we now make the gradient of structural confidence explicit
rather than implicit. The Introduction states that the panel spans ``a gradient of structural
confidence --- PfDHFR carries a co-crystallized inhibitor, PfCRT is represented by a proxy
cavity, and PfClpP and PfATP4 by mechanism-motivated regions without co-crystallized
small-molecule ligands --- so target breadth is read as a pattern of target-specific evidence,
not as four equivalent demonstrations of biological engagement'' (M, \Cref{M-sec:intro}). The
Methods (\Cref{M-sec:methods_docking}) describe each anchor's evidence class, and \Cref{SM-sec:sm_evidence} tabulates the four
evidence classes (catalytic-site ligand, Y01 cavity proxy, catalytic triad, P-type ATPase
machinery). Accordingly, we report all four-target results target by target and never
compute a cross-target composite, so the heterogeneity affects interpretation (breadth as a
hypothesis) rather than a pooled number.

\subsection*{R1.2 --- Validation strategy: feasibility and retrospective validation}
\textbf{Reviewer:} ``The proposed validation strategy lists five categories of experiments but
provides no preliminary data, no estimated feasibility, and no prioritization beyond naming
three candidates. At minimum a retrospective validation against known antimalarial scaffolds
would strengthen the workflow's credibility.''

\textbf{Response:} We strengthened the strategy in three ways.
\begin{enumerate}[leftmargin=1.6em]
\item \textbf{Retrospective validation.} We docked five approved antimalarials
(artemisinin, pyrimethamine, chloroquine, lumefantrine, cipargamin) against the PfDHFR
wild-type/mutant panel (N51I, C59R, S108N, I164L; wild-type receptor stripped of the
co-crystallised MTX to avoid retained-ligand bias, consistent with the mutant
preparation) and the corrected PfCRT panel (LYS-76 WT, K76T, K76A) under the declared
protocols (40 runs, exhaustiveness 32), and computed per-drug RRS profiles to test
whether the framework recovers clinically documented resistance signatures. The result is
honest-negative and is reported as such (\Cref{SM-tab:sm_retrospective_antimalarials}; \Cref{M-sec:disc_validation}): pyrimethamine
RRS is $\approx$100\% across all four PfDHFR mutants (N51I 100.0, C59R 100.0,
S108N 99.7, I164L 100.0), chloroquine RRS is 99.6\% (K76T) and 98.8\%
(K76A) on PfCRT---both within the pipeline-null spread of the corrected
channel (99.4--100.5\%)---and artemisinin/cipargamin show PfDHFR RRS
$\approx$113--118\% while chloroquine/lumefantrine are PfDHFR non-binders
(consistent with their non-antifolate mechanisms). The framework therefore
does not recover the clinical pyrimethamine/chloroquine resistance signatures
under the corrected protocols, which further bounds the interpretation of RRS
as a resistance predictor rather than a biological measurement.
\item \textbf{Prioritization.} \Cref{M-sec:disc_validation} now gives an explicit, ordered five-step experimental
sequence for the five top-priority candidates (PP-15, PP-05, PP-06, PP-11, PP-13):
IC$_{50}$ against wild-type parasites and recombinant proteins $\rightarrow$ isogenic
mutant-line assays (the exact N51I/C59R/S108N/I164L and K76T/K76A panels used here)
$\rightarrow$ orthogonal binding (SPR/ITC) $\rightarrow$ MD refinement $\rightarrow$
ADMET/synthetic-accessibility assessment.
\item \textbf{Feasibility framing.} We now state explicitly that the IC$_{50}$ and mutant-panel
assays are the standard first steps for computational hit triage and require only the five
candidates plus two isogenic panels, and that all claims remain computational hypotheses until
those measurements exist.
\end{enumerate}

\section*{Reviewer 2}

We thank Reviewer 2 for the encouraging framing and for the specific, actionable list. Each
point is addressed below.

\subsection*{R2.1 --- Unreported N\_fav--RRS correlation; power limitation; ``demonstrate''}
\textbf{Reviewer:} Spearman $\rho(N_{\mathrm{fav}},\mathrm{RRS})=+0.515$, $p=0.034$ should be
reported; the PNS--RRS $p=0.020$ ``trend'' and the ``demonstrate'' language are in tension
with \S2.5; with $n=17$ the power to detect $\rho=0.5$ at $\alpha=0.017$ is only 0.37.

\textbf{Response:} The N\_fav--RRS correlation is now reported in \Cref{M-sec:results_descriptors} and \Cref{SM-sec:sm_crossmetric} (Table \Cref{SM-tab:crossmetric}), and was additionally assessed by a 100{,}000-permutation test.
Because the corrected PfCRT re-dock (R2.3) changed the PfCRT column and the
favorability definition (R2.2), the final numbers on the corrected matrix are
$\rho=+0.714$ ($p=0.0013$; permutation $p=0.0019$), significant at the
Bonferroni threshold $\alpha=0.017$. The Discussion now reads the cross-metric
findings with the power limitation stated explicitly: at $n=17$, a two-sided
Spearman test at $\alpha=0.05$ reaches 80\% power only at $|\rho|\ge0.62$, and
at the Bonferroni-corrected $\alpha=0.017$ the power to detect a correlation of
$\rho=0.5$ is only $\approx0.37$ (the reviewer's estimate), so only
moderate-magnitude associations are detectable at this sample size
(\Cref{M-sec:disc_breadth}; \Cref{SM-sec:sm_crossmetric}). The same caveat qualifies the PNS--RRS and
RRS--wild-type-score associations, which are also reported on the corrected
matrix (PNS--RRS $\rho=-0.714$, $p=0.0013$; RRS--mean $|S_{WT}|$ $\rho=+0.691$,
$p=0.0021$; both significant at $\alpha=0.017$). As an additional robustness
check, the companion P2 layer reports that the PNS--RRS association survives and
strengthens after conditioning on molecular weight and Murcko-scaffold prevalence
(partial $\rho=-0.6154$ versus raw $\rho=-0.2098$, $n=12$ complete-two-target
set), consistent with the negative PNS--RRS direction reported here and
indicating that the association is not a size or scaffold artifact.

\subsection*{R2.2 --- The $-6.0$ kcal/mol threshold performs the cross-target comparison}
\textbf{Reviewer:} A single cutoff passes 7/17 (PfDHFR) but 12/17 (PfCRT/PfATP4); favourability
should be defined within each target (rank/percentile) and Table 2 recomputed.

\textbf{Response:} We adopted the reviewer's within-target definition (option b).
Favorability is now defined per target as a Vina estimate at or below that
target's cohort median (PfDHFR $-5.92$, PfCRT $-8.06$, PfClpP $-6.02$, PfATP4
$-6.38$~kcal~mol$^{-1}$ on the corrected matrix; \Cref{SM-tab:sm_vina_distributions}), and \Cref{M-tab:dual_priority} and
the N\_fav--RRS statistics were recomputed accordingly (\Cref{M-sec:methods_pns}, \Cref{M-sec:results_joint}; \Cref{SM-sec:sm_crossmetric}). Under the per-target-median definition, $\rho(N_{\mathrm{fav}},\mathrm{RRS}_{\mathrm{mean}})=+0.714$
($p=0.0013$; permutation $p=0.0019$, 100{,}000 draws), significant at
$\alpha=0.017$. For transparency, the retired single-cutoff variant gave
$\rho=+0.515$ ($p=0.034$); the within-target definition is the one reported
because it avoids a fixed cross-target cutoff, as the reviewer requested.

\subsection*{R2.3 --- PfCRT 6UKJ isoform and grid box (critical)}
\textbf{Reviewer:} 6UKJ is the 7G8 isoform; residue 76 is threonine (K76T); the 28 \AA\ box
excludes residue 76; the PfCRT panel is the sole resilience source for PP-02/PP-06/PP-11/PP-13
and contributes to PP-01/PP-15; anchor to the central negatively charged cavity and construct a
3D7-like wild-type reference.

\textbf{Response:} We agree this was the most serious flaw. We thank the reviewer for the
precise geometric analysis. Three corrections are in place:
\begin{enumerate}[leftmargin=1.6em]
\item \textbf{3D7-like wild-type receptor.} The deposited 6UKJ carries threonine at residue 76;
we restored lysine to obtain a 3D7-like wild-type receptor (LYS-76) and use it as the PfCRT
wild-type reference (\Cref{M-sec:methods_docking}; \Cref{SM-sec:sm_val_crt}).
\item \textbf{Grid anchored to the drug-interaction cavity.} The corrected grid is centered at
$(152.99, 151.042, 159.379)$ with 25 \AA\ dimensions, the V2-corrected box used in our
resistance-resilience analysis. P2Rank pocket prediction independently validates the box: the
top-ranked pocket (score 162.7, probability 0.999) lies 5.7 \AA\ from the grid center (\Cref{SM-sec:sm_val_crt}), i.e., the box covers the central cavity, not a peripheral region.
\item \textbf{Re-dock completed.} All 17 candidates were re-docked against the corrected
LYS-76 receptor (and, for the R2.5 control, the K76T, K76A, and pipeline-null receptors) on
the V2 grid (exhaustiveness 32): 68 runs, all finite. \Cref{SM-tab:vina} (PfCRT column), \Cref{M-tab:rrs_main}
(PfCRT RRS values with the corrected wild-type denominator), and \Cref{M-tab:dual_priority} (N\_fav per-target-
median assignments) were regenerated from these runs. The corrected wild-type scores are
substantially more favorable (mean $-8.42$ vs $-6.19$~kcal~mol$^{-1}$), and the corrected
PfCRT RRS values all fall between 99.1 and 100.6\%. Class counts changed from 6/5/5/1 to
7/4/6/0 (A*/B/C/D), and the corrected RRS range is 73.2--115.6\%.
\end{enumerate}
The corrected-receptor protocol is the one described in the manuscript, and all result tables
were regenerated from the corrected runs (scripts \texttt{p1\_r23\_pfcrt\_redock.py} and
\texttt{p1\_r23\_rrs\_recompute.py} in the repository).

\subsection*{R2.4 --- DEKOIS worse-than-random; retained MTX; Hany et al.}
\textbf{Reviewer:} EF@1\% = EF@5\% = 0.00 suggests retained MTX blocking the antifolate pocket;
re-run on a ligand-stripped receptor; Hany et al. (Drug Des. Devel. Ther. 2025,
10.2147/DDDT.S537065) applied DEKOIS 2.0 to WT and quadruple-mutant PfDHFR with Vina and
rescoring.

\textbf{Response:} We report EF@1\% and EF@5\% (both 0.00) alongside the ROC-AUC in
\Cref{SM-tab:sm_enrichment_validation} and discuss the retained-MTX hypothesis explicitly: the MTX redocking failure
(RMSD $\approx$ 30 \AA) and the PfDHFR grid position indicate that the null enrichment partly
reflects grid positioning rather than scoring-function failure alone (\Cref{M-sec:methods_dock_val}; \Cref{SM-sec:sm_val_dekois}).
We now cite and discuss Hany et al.\ as the closest prior work: they applied the same DEKOIS
2.0 protocol to wild-type and quadruple-mutant PfDHFR with AutoDock Vina and showed that
machine-learning rescoring (CNN-Score, RF-Score-VS) moved the benchmark from worse-than-random
to better-than-random, corroborating our conclusion that enrichment depends on consensus or
rescoring rather than raw Vina ranking (\Cref{SM-sec:sm_val_dekois}).

We performed the requested re-run as a controlled two-arm experiment on the identical DEKOIS
panel (the same 40 actives and 1200 property-matched decoys, same grid): (i) the deposited 7F3Y
receptor with MTX retained, which reproduces the original baseline under our pipeline
(ROC-AUC 0.502 [95\% CI: 0.423--0.586], EF@1\% 0.00, EF@5\% 0.50), and (ii) an MTX-stripped
receptor rebuilt in the same coordinate frame (ROC-AUC 0.563 [95\% CI: 0.477--0.646],
EF@1\% 0.00, EF@5\% 1.00), both at exhaustiveness 32 (the revised-protocol standard, identical
in both arms so the comparison is internally controlled). The result is honest-negative for the
blockade hypothesis: the confidence intervals overlap and EF@1\% remains zero in both arms, so
removing MTX yields only a marginal, non-significant improvement ($\Delta$AUC $= +0.060$) and
does not restore enrichment. We therefore report that the near-chance DEKOIS result is
attributable to the single-method Vina scoring function and grid positioning (near the NADPH
cofactor site) rather than to steric blockade by the retained cofactor ligand, and the manuscript
now states this explicitly with the two-arm data added to \Cref{SM-tab:sm_enrichment_validation}
and \Cref{SM-sec:sm_val_dekois} (\Cref{M-sec:methods_dock_val}; \Cref{M-sec:disc_validation}).

\subsection*{R2.5 --- Batch effect: WT and mutants prepared by different routes}
\textbf{Reviewer:} Provide the required control: WT processed through the mutant pipeline,
use the resulting distribution as a null, place class boundaries outside it; explain the six
RRS values above 100\%.

\textbf{Response:} We now explain the values above 100\% in \Cref{M-sec:results_mutation}: they arise when the
mutant docking estimate is more favorable (more negative) than the wild-type estimate on the
same receptor; they are scoring-function ratios and do not imply biological gain-of-function.
The requested null control was run: the corrected LYS-76 wild-type receptor was passed through
the identical mutation-preparation pipeline without applying any mutation (pipeline-null), and
all 17 candidates were docked against this null receptor and against the K76T and K76A
mutants on the same V2 grid (\Cref{SM-sec:sm_val_redock_null}). The null RRS distribution (null ratios vs the
corrected wild-type denominator) spans 99.4--100.5\% (mean 100.0\%); the K76T ratios span
99.1--100.6\% and K76A 99.4--100.4\%. The maximum deviation of any mutant ratio from its null
ratio is 1.2 percentage points (K76T) and 0.5 (K76A), both far below the 10-point class
boundary spacing. The 80\%/70\% class boundaries therefore fall inside the null spread on the
corrected PfCRT channel, and we do not assert PfCRT-based class distinctions beyond the null
spread: PfCRT RRS values near 100\% are reported as neutral-within-null. The receptor-
preparation contribution to the old PfCRT signal is thus quantified as the entire difference:
the old spread (70.4--110.2\% across K76T/K76A) reflected the raw 6UKJ (7G8, K76T)
receptor used as ``wild type'' and a grid that excluded the drug-interaction cavity.

The neutral-within-null reading is corroborated by two independent companion
analyses in the P2 layer. (i) An external docking replication (39 ligands
$\times$ 8 PfDHFR/PfCRT states, 312/312 finite records, prepared entirely
independently of the present panel) gives a bootstrap RRS distribution of 100.45
[99.59, 101.32], so the corrected PfCRT channel's null spread (99.4--100.5\%)
falls inside the pipeline noise of a fully separate preparation. (ii) In the
companion P2 layer, RRS class assignments on its canonical panel are robust to
score perturbation: under bounded $\pm1$ kcal/mol perturbations of individual
docking scores (272 candidate--score records), the class was unchanged in
212/272 (77.9\%), and leave-one-mutant-out and wild-type-threshold (4 vs
5 kcal/mol) sensitivity runs leave its reported class distributions stable. Separately, an MD-based MM-GBSA $\Delta\Delta G$ analysis
in the companion layer found 0 of 12 PfDHFR/PfCRT mutants significantly weaker
and 1 significantly tighter at 95\% CI (max $|\Delta\Delta G|=5.37$ kcal/mol),
corroborating that mutant-to-wild-type ratios above 100\% reflect retention
rather than gain.

\subsection*{R2.6 --- Repository returns 404; data unavailable}
\textbf{Reviewer:} The repository URL returns 404; the wild-type denominators do not appear
elsewhere; provide a maintained repository with license, release tag, and archival DOI.

\textbf{Response:} [AUTHOR ACTION PENDING --- make
\url{https://github.com/NanaEngo/Malaria_codesV2} public, add a software license and a
versioned release tag, deposit the release on Zenodo, and update the Data Availability
statement (M) and \Cref{SM-sec:sm_repro} with the DOI.] In parallel, the WT/mutant score table required to
recompute \Cref{M-tab:rrs_main} is now included directly in the SI (\Cref{SM-tab:sm_wt_mutant_scores}: the
PfDHFR wild-type/mutant panel and the corrected PfCRT LYS-76/K76T/K76A scores underlying
every RRS denominator), so the RRS denominators no longer reside only in the repository.

\subsection*{R2.7 --- PNS and ACSI never defined}
\textbf{Response:} \Cref{M-sec:methods_pns} now gives the formal definitions. ACSI =
$0.40\times D_{\mathrm{DrugBank}} + 0.25\times D_{\mathrm{ANPDB}} + 0.20\times f_{sp^3}
+ 0.15\times \mathrm{NPL}$ with min--max normalization over the 17-candidate cohort
($D$ = minimum Tanimoto distance to approved-drug / African-NP reference sets, $f_{sp^3}$ =
fraction of sp$^3$ carbons, NPL = natural-product-likeness). PNS is computed from a STRING
protein--protein interaction network at confidence threshold 700, with PfCRT centrality imputed
as the network mean (0.151); \Cref{SM-tab:s_pns_sensitivity} reports the imputation sensitivity (Spearman
$\rho$ = 0.963--1.000 across zero to double the imputed value). No PNS or ACSI value is
interpreted causally.

\subsection*{R2.8 --- \Cref{SM-tab:sm_redocking_validation} vs \Cref{SM-sec:sm_val_redock} contradiction}
\textbf{Reviewer:} \S12.2 says no pose-reproduction rate is claimed for PfClpP, yet Table S8
lists ``1 success'' for all four targets; 9N10 has no small-molecule ligand; the score-stratified
MMV rows are circular.

\textbf{Response:} The validation tables were rebuilt to remove the contradiction.
\Cref{SM-tab:sm_redocking_validation,SM-tab:sm_validation_datasets} now separate (i) full-ligand redocking (PfDHFR P218, RMSD 1.42 \AA; PfCRT Y01 proxy,
1.78 \AA) from (ii) fragment pose reproduction (PfClpP peptide, 1.65 \AA) and cofactor pose
reproduction (PfATP4 ADP, 1.23 \AA), and state explicitly that (ii) is not equivalent to
full-ligand redocking and that 9N10 contains no co-crystallized small-molecule ligand. The MMV
rows for PfCRT/PfATP4 are explicitly labelled score-stratified retrodictive consistency checks,
not external discrimination. The text in \Cref{SM-sec:sm_val_redock} was revised to match.

\subsection*{R2.9 --- Single seed; PP-15 marginally above threshold}
\textbf{Reviewer:} PP-15 exceeds the cutoff by 0.045 (PfDHFR) and 0.084 (PfCRT); run at least
five seeds and report dispersion.

\textbf{Response:} Five-seed sensitivity runs (seeds 101--505) are reported in \Cref{SM-tab:sm_multiseed_validation}
for PP-15 and PP-01 on PfDHFR and PfCRT; mode-1 affinities are stable to seed choice
(maximum spread 0.054 kcal/mol). The table transparently states that those runs used the
P2 canonical-grid protocol and therefore report seed dispersion for that protocol. After the
R2.3 re-dock on the corrected receptor and declared grids, the per-target-median
favorability assignments (R2.2) place 48 of 68 candidate--target pairs at least
0.3 kcal mol$^{-1}$ from the relevant target median; the remaining 20 pairs lie within
0.3 kcal mol$^{-1}$, including four that sit exactly at the median (PP-05 PfATP4, PP-10
PfClpP, PP-12 PfCRT, PP-13 PfDHFR), so the within-target favorability counts should be read
with this boundary sensitivity in mind. The corrected-protocol single-run scores used for those assignments
are the same values reported in the corrected \Cref{SM-tab:vina}, and the re-dock used a single seed
with exhaustiveness 32 as in the original protocol. The seed dispersion shown in \Cref{SM-tab:sm_multiseed_validation}
($<0.06$ kcal mol$^{-1}$) was measured under the P2 canonical-grid protocol and therefore
provides a scale reference for seed sensitivity rather than a direct bound on the corrected
single-seed V2-grid runs; the boundary analysis above (20 pairs within 0.3 kcal mol$^{-1}$ of a
target median, four exactly at the median) identifies which N\_fav assignments are most
sensitive to single-seed variation.

\subsection*{R2m --- Minor points}
\begin{enumerate}[leftmargin=1.6em]
\item \textbf{Duplicated paragraph (pp. 17--18):} the verbatim duplicate is removed; the
top-candidate summary now appears once in \Cref{M-sec:results_joint} and is not repeated verbatim in the
Discussion.
\item \textbf{VAE citation:} we now cite G\'omez-Bombarelli et al.\ (ACS Cent. Sci.\ 2018)
for the SMILES VAE (\Cref{M-sec:methods_funnel}).
\item \textbf{``99.3\% reduction'':} reframed as an absolute reduction (263{,}424 $\rightarrow$
1{,}936 calculations) with the explicit caveat that this describes computational cost, not
active-compound retention, and that a retention experiment using known actives would be needed
(\Cref{M-sec:disc_efficiency}).
\item \textbf{RRS eligibility notation:} all instances now use $|S_{\mathrm{Vina},i,\mathrm{WT},t}|$
consistently (\Cref{M-sec:methods_rrs}; \Cref{SM-sec:sm_rrs}); no $\Delta G$ notation remains.
\item \textbf{Section S12 numbered twice; decoy counts:} SI now contains a single S12 block
(S12.1--S12.7); the DEKOIS decoy count is consistently 1{,}200.
\item \textbf{Figure 2 annotations:} the figure was regenerated with $\rho$, $p$, and $n$
annotated in each panel from the corrected canonical data (PNS--RRS $\rho=-0.714$,
$p=0.0013$; ACSI--RRS $\rho=-0.190$, $p=0.465$; $n=17$ in both panels).
\item \textbf{MPO sensitivity deferred:} the essential results are now summarized in \Cref{SM-sec:sm_mpo_weights}
(25 weight perturbations; Spearman $\rho=0.792$ for the top-1000 ranking), so the manuscript
stands alone while the full methodology remains in the companion chemical-space manuscript.
\end{enumerate}

\section*{Status of items formerly marked [TO COMPLETE]}
Completed for this revision: (i) corrected-receptor PfCRT re-dock of the 17-candidate panel
and regeneration of \Cref{SM-tab:vina}/\Cref{M-tab:rrs_main}/\Cref{M-tab:dual_priority} (R2.3/R2.9); (ii) the batch-effect null
control and class-boundary placement (R2.5; \Cref{SM-sec:sm_val_redock_null}); (iii) Figure 2 statistical
annotations (R2m.6); (iv) per-target-median favorability and recomputed
\Cref{M-tab:dual_priority}/N\_fav--RRS (R2.2); (v) retrospective validation against five approved
antimalarials (R1.2; 40 runs, honest-negative result reported in \Cref{SM-tab:sm_retrospective_antimalarials} and
\Cref{M-sec:disc_validation}); (vi) the WT/mutant score table in the SI (R2.6;
\Cref{SM-tab:sm_wt_mutant_scores}); and (vii) the controlled two-arm MTX-retained versus
MTX-stripped DEKOIS re-run (R2.4; honest-negative result, two rows added to
\Cref{SM-tab:sm_enrichment_validation} and discussed in \Cref{SM-sec:sm_val_dekois}).
Remaining, pending author action: repository release + Zenodo DOI (R2.6; checklist provided in
the repository).

\vspace{2em}
\noindent Sincerely,\\
Myke Vital Sao Temgoua, Jean-Pierre Tchapet Njafa, Penabei Samafou,\\
Wilfred Fon Mbacham, Serge Guy Nana Engo

\end{document}